\chapter{Theoretical Background}

\section{Priority Queues}

A priority queue stores items under an ordering relation and supports repeated access to the item with highest priority.
In heapx, the highest-priority item is the one that compares first according to the user-supplied comparator.
The API names this item the minimum because the natural comparator for ascending numeric keys returns smaller values first.

Priority queues are frequently analyzed through the cost of insertion, minimum extraction, and priority decrease.
The relative importance of these operations depends on the workload.
For example, Dijkstra's shortest-path algorithm performs extract-min operations and may perform many decrease-key operations \cite{dijkstra1959note}.

\section{Binary Heap}

The binary heap stores elements in an array and interprets the array as a complete binary tree \cite{cormen2009introduction}.
For a zero-based array index $i$, the structural relationships are:

\begin{align}
    parent(i) & = \left\lfloor \frac{i - 1}{2} \right\rfloor, \\
    left(i)   & = 2i + 1,                                     \\
    right(i)  & = 2i + 2.
\end{align}

The min-heap invariant requires every parent to compare no worse than its children.
Insertion appends a new item at the end of the array and restores order by sifting upward.
Extracting the minimum removes the root, moves the last entry to the root, and restores order by sifting downward.

The binary heap is simple and cache-friendly.
Its main costs are:

\begin{itemize}
    \item insert: $O(\log n)$;
    \item decrease-key: $O(\log n)$ when the item position is known;
    \item remove: $O(\log n)$ when the item position is known;
    \item peek-min: $O(1)$;
    \item extract-min: $O(\log n)$.
\end{itemize}

\section{Fibonacci Heap}

Fibonacci heaps use a collection of heap-ordered trees and obtain strong amortized bounds \cite{fredman1987fibonacci}.
Insertion adds a singleton tree to the root list.
Decrease-key cuts a node from its parent if heap order is violated and then applies cascading cuts.
Extract-min removes the minimum root, promotes its children, and consolidates roots so that no two roots have the same degree.

The relevant amortized bounds are:

\begin{itemize}
    \item insert: amortized $O(1)$;
    \item decrease-key: amortized $O(1)$;
    \item remove: amortized $O(\log n)$;
    \item peek-min: $O(1)$;
    \item extract-min: amortized $O(\log n)$.
\end{itemize}

The implementation cost is higher than for a binary heap.
Nodes need parent, child, sibling links, degree metadata, and mark bits.
This is a useful trade-off to study because theoretical complexity and measured performance can diverge when pointer chasing, allocation, and cache locality dominate.

\section{Simple Fibonacci or Kaplan Heap}

The Kaplan backend follows the simple Fibonacci heap idea described by Kaplan, Tarjan, and Zwick \cite{kaplan2014fibonacci}.
In the heapx implementation, the steady-state representation is a single heap-ordered tree.
Insertion links a singleton node against the root using a naive link.
Extract-min removes the root, treats its children as temporary roots, performs fair links among roots with equal rank, and then folds the remaining roots back into one tree.

Decrease-key differs from the classic Fibonacci heap.
Instead of cascading cuts, the Kaplan backend performs cascading rank and state changes followed by one structural cut.
This makes it an interesting contrast with the classic Fibonacci backend: both aim at efficient decrease-key behavior, but their internal repairs are different.

\section{Worst-Case and Amortized Complexity}

The report distinguishes worst-case bounds from amortized bounds.
The binary heap operations have straightforward worst-case logarithmic behavior.
The Fibonacci-style heaps rely on amortized analysis: individual operations may perform more work, but sequences of operations satisfy the documented average-over-sequence bounds.

This distinction matters when interpreting benchmarks.
A benchmark result is an empirical measurement of this implementation, compiled with a particular toolchain, running on a particular machine, under a particular workload.
It should be read together with the theoretical bounds, not as a substitute for them.
